\begin{theorem}[Determination by the seven native transports]
    \label{thm:peps_two_plaquette_transport_table}
    \lean{TNLean.PEPS.twoPlaquette_internal_eq_of_edgeTransports,TNLean.PEPS.verticalTwoPlaquette_internal_eq_of_transports,TNLean.PEPS.translatedTwoPlaquette_internal_eq_of_transports,TNLean.PEPS.torusVerticalFluxAssignment_up_transport,TNLean.PEPS.torusVerticalFluxAssignment_right_transport_of_ne}
    \leanok
    Suppose an induced region is isomorphic to the six-vertex two-plaquette
    graph. Equality of the directed transports on its seven mapped edges
    implies equality of every internal ordered bond coefficient. On a native
    torus this applies to a translated horizontal tile when its periods satisfy
    $w\ge4$ and $H\ge3$, and to a vertical tile when $w\ge3$ and $H\ge4$.
    For the vertical movement assignment, every upward transport is the
    identity, as is every rightward transport outside its two selected rows.
    These finite-block statements are auxiliary to
    \cite[Theorem~6.16]{Schuch2010PEPS}.
\end{theorem}
\begin{proof}
    \leanok
    \uses{thm:peps_two_plaquette_graph,thm:peps_translated_two_plaquette_geometry,thm:peps_vertical_two_plaquette_geometry,thm:peps_internal_gauge_transport}
    The graph isomorphism accounts for every internal edge. Each directed
    equality determines the corresponding ordered coefficient. The native
    forms follow by evaluating the vertex coordinates and reversing the two
    necessary traversals. The movement assignment is supported on two
    horizontal bonds, which proves the remaining transport statements.
\end{proof}
